\section{Multiplicities and the Cartan Matrix}

\begin{definition}[Multiplicity of a simple module]
    Let \(A\) be an Artinian ring, and let \(V\) be a finitely generated left \(A\)-module.
    Then \(V\) has a composition series
    \[
        0=V_n \subset V_{n-1} \subset \cdots \subset V_1 \subset V_0=V,
    \]
    where each quotient
    \[
        V_{i-1}/V_i
    \]
    is a simple \(A\)-module.

    For a simple \(A\)-module \(S\), the number of composition factors
    \[
        V_{i-1}/V_i
    \]
    which are isomorphic to \(S\) is called the \emph{multiplicity of \(S\) in \(V\)}.
    We denote it by
    \[
        [V:S].
    \]
\end{definition}

\begin{theorem}
    Let \(A\) be an Artinian ring, and let \(f\in A\) be a primitive idempotent. Put
    \[
        J=J(A),\qquad \overline{A}=A/J,
        \qquad e=f+J\in \overline{A}.
    \]
    Assume that \(e\) is a primitive idempotent in \(\overline{A}\). Then, for every finitely
    generated left \(A\)-module \(V\), the multiplicity of the simple \(A\)-module
    \(\overline{A}e\) in \(V\) is equal to the composition length of \(fV\) as a left
    \(fAf\)-module. In other words,
    \[
        [V:\overline{A}e]=\ell_{fAf}(fV).
    \]
\end{theorem}

\begin{proof}
    Let
    \[
        V=V_0\supset V_1\supset \cdots \supset V_n=0
    \]
    be a composition series of \(V\). Applying \(f\) gives a descending chain of
    left \(fAf\)-modules
    \[
        fV=fV_0\supseteq fV_1\supseteq \cdots \supseteq fV_n=0.
    \]

    For each \(i\), the quotient \(V_i/V_{i+1}\) is a simple \(A\)-module, hence it is also
    a simple \(\overline{A}\)-module because \(J(A)\) annihilates every simple
    \(A\)-module. We claim that
    \[
        V_i/V_{i+1}\cong \overline{A}e
        \quad\Longleftrightarrow\quad
        f(V_i/V_{i+1})\neq 0.
    \]
    Indeed, the action of \(f\) on a simple \(A\)-module is the same as the action of
    \(e=f+J\). Hence, for a simple \(A\)-module \(S\),
    \[
        fS\neq 0
        \quad\Longleftrightarrow\quad
        eS\neq 0.
    \]
    Moreover,
    \[
        eS\cong \operatorname{Hom}_{\overline{A}}(\overline{A}e,S),
    \]
    and since \(\overline{A}e\) and \(S\) are simple \(\overline{A}\)-modules, this space is
    nonzero if and only if
    \[
        S\cong \overline{A}e.
    \]
    Thus the claim follows.

    Now for each \(i\), we have
    \[
        f(V_i/V_{i+1})
        \cong
        \frac{fV_i+V_{i+1}}{V_{i+1}}
        \cong
        \frac{fV_i}{fV_i\cap V_{i+1}}.
    \]
    Since \(f\) is idempotent, one has
    \[
        fV_i\cap V_{i+1}=fV_{i+1}.
    \]
    Indeed, if \(x\in fV_i\cap V_{i+1}\), then \(x=fv\) for some \(v\in V_i\), and since
    \(x\in V_{i+1}\), we get
    \[
        x=fx\in fV_{i+1}.
    \]
    The reverse inclusion is clear. Therefore
    \[
        f(V_i/V_{i+1})
        \cong
        \frac{fV_i}{fV_{i+1}}.
    \]

    Hence every quotient
    \[
        fV_i/fV_{i+1}
    \]
    is either \(0\) or a simple \(fAf\)-module. Moreover, it is nonzero precisely when
    \[
        V_i/V_{i+1}\cong \overline{A}e.
    \]
    Therefore, after deleting the zero quotients from the chain
    \[
        fV=fV_0\supseteq fV_1\supseteq \cdots \supseteq fV_n=0,
    \]
    we obtain a composition series of \(fV\) as an \(fAf\)-module.

    Consequently, the composition length of \(fV\) as an \(fAf\)-module is exactly the
    number of composition factors of \(V\) isomorphic to \(\overline{A}e\). Hence
    \[
        \ell_{fAf}(fV)=[V:\overline{A}e].
    \]
\end{proof}

\begin{definition}[Cartan invariant and Cartan matrix]
    Let \(A\) be an Artinian ring, and let
    \[
        J=J(A), \qquad \overline{A}=A/J.
    \]
    Let
    \[
        1=e_1+\cdots+e_n
    \]
    be a primitive idempotent decomposition of \(1\) in \(\overline{A}\). Choose a lifting
    \[
        1=f_1+\cdots+f_n
    \]
    to a primitive idempotent decomposition of \(1\) in \(A\), so that
    \[
        e_i=f_i+J(A)
        \qquad
        \text{for all } 1\le i\le n.
    \]

    For \(1\le i,j\le n\), define
    \[
        c_{ij}:=\bigl[Af_j:\overline{A}e_i\bigr],
    \]
    that is, \(c_{ij}\) is the multiplicity of the simple \(A\)-module
    \(\overline{A}e_i\) as a composition factor of the \(A\)-module \(Af_j\).

    The integer \(c_{ij}\) is called the \emph{Cartan invariant} of \(A\). The matrix
    \[
        C_A=(c_{ij})_{1\le i,j\le n}
    \]
    is called the \emph{Cartan matrix} of \(A\).
\end{definition}

\begin{theorem}
    Let \(A\) be an Artinian ring, and let
    \[
        J=J(A), \qquad \overline{A}=A/J.
    \]
    Let
    \[
        1=e_1+\cdots+e_n
    \]
    be a primitive idempotent decomposition in \(\overline{A}\), and let
    \[
        1=f_1+\cdots+f_n
    \]
    be a primitive idempotent decomposition in \(A\) such that
    \[
        e_i=f_i+J(A)
        \qquad
        \text{for all } 1\le i\le n.
    \]
    Let
    \[
        c_{ij}=\bigl[Af_j:\overline{A}e_i\bigr]
    \]
    be the Cartan invariant of \(A\). Then
    \[
        c_{ij}\neq 0
        \quad\Longleftrightarrow\quad
        f_iAf_j\neq 0.
    \]
\end{theorem}

\begin{proof}
    By the multiplicity theorem for primitive idempotents, for any finitely generated
    \(A\)-module \(V\), the multiplicity of \(\overline{A}e_i\) in \(V\) is equal to the
    composition length of \(f_iV\) as an \(f_iAf_i\)-module. That is,
    \[
        [V:\overline{A}e_i]
        =
        \ell_{f_iAf_i}(f_iV).
    \]

    Now take
    \[
        V=Af_j.
    \]
    Then
    \[
        f_iV=f_i(Af_j)=f_iAf_j.
    \]
    Therefore
    \[
        c_{ij}
            =
            [Af_j:\overline{A}e_i]
        =
        \ell_{f_iAf_i}(f_iAf_j).
    \]

    Since \(f_iAf_j\) is a finite length \(f_iAf_i\)-module, its composition length is nonzero
    if and only if the module itself is nonzero. Hence
    \[
        c_{ij}\neq 0
        \quad\Longleftrightarrow\quad
        \ell_{f_iAf_i}(f_iAf_j)\neq 0
        \quad\Longleftrightarrow\quad
        f_iAf_j\neq 0.
    \]
    This proves the theorem.
\end{proof}

\begin{proposition}
    Let \(A\) be a finite-dimensional \(k\)-algebra over a field \(k\). Let \(S\) be a
    simple left \(A\)-module with projective cover
    \[
        P_S \twoheadrightarrow S,
    \]
    and let \(U\) be a finite-dimensional left \(A\)-module.

    \begin{enumerate}
        \item If \(T\) is a simple left \(A\)-module, then
              \[
                  \dim_k \operatorname{Hom}_A(P_S,T)
                  =
                  \begin{cases}
                      \dim_k \operatorname{End}_A(S), & S\cong T,     \\
                      0,                              & S\not\cong T.
                  \end{cases}
              \]

        \item The multiplicity of \(S\) in \(U\) is given by
              \[
                  [U:S]
                  =
                  \frac{\dim_k \operatorname{Hom}_A(P_S,U)}
                  {\dim_k \operatorname{End}_A(S)}.
              \]

        \item If \(e\in A\) is an idempotent, then
              \[
                  \dim_k \operatorname{Hom}_A(Ae,U)=\dim_k eU.
              \]

        \item Let \(e_S,e_T\in A\) be idempotents such that
              \[
                  P_S=Ae_S,
                  \qquad
                  P_T=Ae_T
              \]
              are projective covers of \(S\) and \(T\), respectively. If
              \[
                  c_{ST}:=[P_T:S]
              \]
              denotes the Cartan invariant, then
              \[
                  c_{ST}
                  =
                  \frac{\dim_k \operatorname{Hom}_A(P_S,P_T)}
                  {\dim_k \operatorname{End}_A(S)}
                  =
                  \frac{\dim_k e_SAe_T}
                  {\dim_k \operatorname{End}_A(S)}.
              \]
              In particular, if \(k\) is algebraically closed, then
              \[
                  c_{ST}
                  =
                  \dim_k \operatorname{Hom}_A(P_S,P_T)
                  =
                  \dim_k e_SAe_T.
              \]
    \end{enumerate}
\end{proposition}

\begin{proof}
    Since \(P_S\twoheadrightarrow S\) is a projective cover, we have
    \[
        P_S/\operatorname{Rad}(P_S)\cong S.
    \]

    \medskip

    \noindent
    \textbf{(1)}
    Let \(T\) be a simple left \(A\)-module. Every homomorphism
    \[
        \varphi:P_S\to T
    \]
    vanishes on \(\operatorname{Rad}(P_S)\). Indeed, if \(\varphi=0\), this is clear. If
    \(\varphi\neq 0\), then \(\varphi\) is surjective because \(T\) is simple, hence
    \(\ker\varphi\) is a maximal submodule of \(P_S\). Therefore
    \[
        \operatorname{Rad}(P_S)\subseteq \ker\varphi.
    \]
    Thus every map \(P_S\to T\) factors uniquely through
    \(P_S/\operatorname{Rad}(P_S)\). Hence
    \[
        \operatorname{Hom}_A(P_S,T)
        \cong
        \operatorname{Hom}_A(P_S/\operatorname{Rad}(P_S),T)
        \cong
        \operatorname{Hom}_A(S,T).
    \]
    By Schur's lemma,
    \[
        \operatorname{Hom}_A(S,T)=0
    \]
    if \(S\not\cong T\), while
    \[
        \operatorname{Hom}_A(S,S)=\operatorname{End}_A(S)
    \]
    if \(S\cong T\). Therefore
    \[
        \dim_k \operatorname{Hom}_A(P_S,T)
        =
        \begin{cases}
            \dim_k \operatorname{End}_A(S), & S\cong T,     \\
            0,                              & S\not\cong T.
        \end{cases}
    \]

    \medskip

    \noindent
    \textbf{(2)}
    Let
    \[
        U=U_0\supset U_1\supset \cdots \supset U_m=0
    \]
    be a composition series of \(U\). Since \(P_S\) is projective, the functor
    \[
        \operatorname{Hom}_A(P_S,-)
    \]
    is exact. Therefore, for each \(i\), applying \(\operatorname{Hom}_A(P_S,-)\) to
    \[
        0\longrightarrow U_{i+1}\longrightarrow U_i
        \longrightarrow U_i/U_{i+1}\longrightarrow 0
    \]
    gives an exact sequence
    \[
        0\longrightarrow
        \operatorname{Hom}_A(P_S,U_{i+1})
        \longrightarrow
        \operatorname{Hom}_A(P_S,U_i)
        \longrightarrow
        \operatorname{Hom}_A(P_S,U_i/U_{i+1})
        \longrightarrow 0.
    \]
    Hence
    \[
        \dim_k\operatorname{Hom}_A(P_S,U)
        =
        \sum_{i=0}^{m-1}
        \dim_k\operatorname{Hom}_A(P_S,U_i/U_{i+1}).
    \]
    By part \((1)\), each composition factor \(U_i/U_{i+1}\) contributes
    \(\dim_k\operatorname{End}_A(S)\) precisely when
    \[
        U_i/U_{i+1}\cong S,
    \]
    and contributes \(0\) otherwise. Therefore
    \[
        \dim_k\operatorname{Hom}_A(P_S,U)
        =
        [U:S]\cdot \dim_k\operatorname{End}_A(S).
    \]
    Thus
    \[
        [U:S]
        =
        \frac{\dim_k \operatorname{Hom}_A(P_S,U)}
        {\dim_k \operatorname{End}_A(S)}.
    \]

    \medskip

    \noindent
    \textbf{(3)}
    Define a map
    \[
        \Phi:\operatorname{Hom}_A(Ae,U)\longrightarrow eU
    \]
    by
    \[
        \Phi(\varphi)=\varphi(e).
    \]
    Since
    \[
        \varphi(e)=\varphi(e^2)=e\varphi(e),
    \]
    we indeed have \(\varphi(e)\in eU\).

    Conversely, for \(u\in eU\), define
    \[
        \psi_u:Ae\longrightarrow U
    \]
    by
    \[
        \psi_u(ae)=au.
    \]
    This is well-defined and \(A\)-linear. Moreover,
    \[
        \Phi(\psi_u)=\psi_u(e)=u.
    \]
    Also, for \(\varphi\in\operatorname{Hom}_A(Ae,U)\),
    \[
        \psi_{\varphi(e)}(ae)=a\varphi(e)=\varphi(ae).
    \]
    Hence \(\Phi\) is an isomorphism of \(k\)-vector spaces:
    \[
        \operatorname{Hom}_A(Ae,U)\cong eU.
    \]
    Therefore
    \[
        \dim_k \operatorname{Hom}_A(Ae,U)=\dim_k eU.
    \]

    \medskip

    \noindent
    \textbf{(4)}
    By definition,
    \[
        c_{ST}=[P_T:S].
    \]
    Applying part \((2)\) to \(U=P_T\), we obtain
    \[
        c_{ST}
        =
        \frac{\dim_k\operatorname{Hom}_A(P_S,P_T)}
        {\dim_k\operatorname{End}_A(S)}.
    \]
    Since \(P_S=Ae_S\) and \(P_T=Ae_T\), part \((3)\) gives
    \[
        \operatorname{Hom}_A(P_S,P_T)
        =
        \operatorname{Hom}_A(Ae_S,Ae_T)
        \cong
        e_SAe_T.
    \]
    Thus
    \[
        \dim_k\operatorname{Hom}_A(P_S,P_T)
        =
        \dim_k e_SAe_T.
    \]
    Therefore
    \[
        c_{ST}
        =
        \frac{\dim_k \operatorname{Hom}_A(P_S,P_T)}
        {\dim_k \operatorname{End}_A(S)}
        =
        \frac{\dim_k e_SAe_T}
        {\dim_k \operatorname{End}_A(S)}.
    \]

    If \(k\) is algebraically closed, then by Schur's lemma
    \[
        \operatorname{End}_A(S)\cong k,
    \]
    so
    \[
        \dim_k\operatorname{End}_A(S)=1.
    \]
    Hence
    \[
        c_{ST}
        =
        \dim_k \operatorname{Hom}_A(P_S,P_T)
        =
        \dim_k e_SAe_T.
    \]
\end{proof}

